\SourceLeaf{183}selbst der vollendeten Unentschiedenheit in dieser Sache, dann, wenn jener
Beweis gegen die Euklideische Geometrie, wie ich nicht erwarten durfte, nicht
für stringent zu halten sei --- [sei von mir erkannt worden]; theils, weil doch
alle die vielen bisherigen Untersuchungen aus der ebenen Geometrie nicht für
verloren zu achten: sondern mit wenigen Modificationen zu gebrauchen wären,
und denn doch wenigstens bis zu einer ziemlich weiten Grenze hin, die sich
vielleicht noch näher bestimmen liesse, auch die Sätze der körperlichen
Geometrie und der Mechanik näherungsweise Gültigkeit hätten; fand ich heute
Abend --- eben mit einem Versuch beschäftigt, einen Eingang zu Ihrer
transcendenten Trigonometrie zu finden, und weil es mir nicht gelingen
wollte, in der Ebene dafür hinreichende, bestimmte Functionen zu erhalten,
zu räumlichen Constructionen fortgehend, zu meiner nicht geringen
Überraschung folgenden \textls{Beweis für die Euklideische Parallel-Theorie}. .\ .\ .\ .\ .

.\ .\ .\ .\ .\ .\ Gerade im Begriff zu schliessen bemerke ich noch, dass der obige
Beweis für die Euklideische Parallel-Theorie fehlerhaft ist. .\ .\ .\ .\ . Also wäre
auch hier die Hoffnung verschwunden, zu einem völlig entschiedenen Resultat
zu kommen, und ich muss mich wieder bei dem vorhin Angeführten beruhigen.
Übrigens glaube ich auf jenem Wege wenigstens einen Schritt zu Ihrer
transcendenten Trigonometrie gethan zu haben, indem ich, mit Hülfe der
sphärischen Trigonometrie, die Verhältnisse aller Constanten, wenigstens
\textls{durch Construction der rechtwinkligen} Dreiecke angeben kann. Es fehlt mir
noch die wirkliche Berechnung der Basis eines gleichschenkligen Dreiecks aus
der Seite, wofür ich suchen werde vom gleichseitigen Dreieck auszugehen.
.\ .\ .\}

\begin{center}\rule{3cm}{0.4pt}\end{center}

\Needspace{11\baselineskip}
\begin{center}\textsc{Bemerkung.}\end{center}

In dem Briefe vom 12. December 1816, von dem hier Stücke mitgetheilt sind,
bezieht sich \textsc{Wachter} wiederholt auf ein Gespräch mit \textsc{Gauss},
das während seines »letzten Aufenthalts in Göttingen« stattgefunden und die
antieuklidische Geometrie zum Gegenstand gehabt hatte. Aus einem Briefe von
\textsc{Gauss} an \textsc{Olbers} vom 4. Juni 1816 geht hervor, dass dieser
Besuch \textsc{Wachters} in den April 1816 zu setzen ist.

\begin{flushright}\textsc{Stäckel.}\end{flushright}

\begin{center}\rule{3cm}{0.4pt}\end{center}

